% 第一章：项目介绍

\section{项目介绍} \label{sec:intro}

\subsection{问题背景}

多核处理器上，任务的执行时间不仅取决于指令条数，还取决于它运行时数据在缓存层次中的位置。对于有向无环图（DAG）任务，同一作业的不同节点往往复用同一批输入数据；若这些节点被分散到不同核心，每个核心都要重新把数据从末级缓存或内存搬入私有缓存，原本可以命中的访问退化为长延迟访问。调度器在"把任务放在哪个核心"上做的选择，因此直接决定了缓存重用能否发生。

传统做法把这个问题拆成两半：实时调度文献关注截止期与可调度性，通常假设执行时间与放置无关；缓存感知研究关注重用距离与缓存冲突，但往往在离线或固定分配的前提下建模。Yat-CAShed 的出发点是把两者接起来：用缓存重用距离估计"这次放置能省多少时间"，再把这个估计值接进在线选核决策，并落到真实的 Linux 调度类里。

\subsection{项目目标与交付物}

本项目要做的是可运行、可复现的系统，而不是纸面模型。主交付物是内核侧的一套东西，见\autoref{tab:deliverables}。

\begin{table}[H]
\centering
\caption{交付物与规模}
\label{tab:deliverables}
\small
\begin{tabular}{@{}lrl@{}}
\toprule
载体 & 规模 & 入口 \\
\midrule
\texttt{kernel/sched/yat\_casched.c} & 2752 行 C & \texttt{make kernel-build} \\
\texttt{include/linux/sched/yat\_casched.h} & 266 行 & \texttt{make symbol-check} \\
\texttt{include/trace/events/yat\_casched.h} & 228 行 & \texttt{debugfs} 寄存器组 \\
\texttt{yat\_test\_scripts/} & 3 配置 $\times$ 5 次 & \texttt{compare\_cfs\_vs\_yat.sh} \\
\texttt{patches/}（7 个补丁） & —— & \texttt{make patch-check} \\
\bottomrule
\end{tabular}
\end{table}

表中没有仿真器。作者在项目前期写过一个 Java 模拟器（\texttt{simulator/}）用于设计期校验，但它的执行时间与命中判定都是理想化模型，与内核不同构，不能作为"算法在真实系统上有效"的证据。本轮复核后已把它从证据链中移除，正文只在第\ref{sec:test}节用一小段交代其结论与去处，全部数据留在仓库文档里。

\subsection{本文的主要工作}

本文不是对已有结论的复述，而是把系统做完、测完并核对之后写成的工程记录。主要工作有四件：

\begin{enumerate}
    \item \textbf{实现}：在 Linux 6.8 上新增内置调度类 \texttt{SCHED\_YAT\_CASCHED}（策略号 8），完成选核、DAG/作业生命周期、每核历史表、加速表、28 项计数器与 \texttt{debugfs} 观测接口，并给出可独立应用的补丁包。
    \item \textbf{修复}：定位并修掉一处被 \texttt{lockdep} 实证的锁反序（\texttt{task\_dead} 在开中断上下文销毁作业，与 balance 路径构成 ABBA），重建镜像并复跑对照，同类告警计数归零。
    \item \textbf{实测}：在 QEMU（4 vCPU，TCG）上用同一份负载跑三组对照、每组重复 5 次，给出中位数与全距；用计数器证明调度类回调确实被调用、决策确实落地（接管率 $98.5\%$），并用"单核 pin"与"无唤醒"两组对照把墙钟代价定位到唤醒路径的回调，而不是 AJLR 选核本身。
    \item \textbf{可复现}：每个数字都能由仓库内脚本重算——\texttt{make verify-all} 是静态门禁，\texttt{scripts/run\_kernel\_repeats.sh} 是实测脚本，结果按配置归档在 \texttt{docs/validation/}。
\end{enumerate}

\subsection{阅读指引}

第\ref{sec:realtime}节交代背景与相关工作；第\ref{sec:design}节给出架构与模块边界；第\ref{sec:theory}节给出算法与它到内核实现的映射；第\ref{sec:kernel}节讲内核实现细节与锁序；第\ref{sec:test}节给出实测协议、计数器证据与代价区间；第\ref{sec:difficulties}节记录工程问题与解决过程；第\ref{sec:conclusion}节总结并给出后续改进路线。